\documentclass[10pt,a4paper]{amsart}
\usepackage[margin=19mm]{geometry}
\usepackage{amsmath,amssymb,mathtools}
\usepackage[scr=rsfs,cal=euler]{mathalfa}
\usepackage{xfrac}
\usepackage[unicode,hidelinks,bookmarks=false]{hyperref}
\newcommand{\ie}{i.e.\ }
\newcommand{\cf}{cf.\ }
\newcommand{\resp}{respectively\ }
\newcommand{\Subsection}{\subsection}
\providecommand{\nobreakdash}{\nobreak\hbox{-}}
\setlength{\emergencystretch}{2em}
\multlinegap0pt
\usepackage{mathtools}
\usepackage{amssymb}
\usepackage{tikz}
\usepackage{xfrac}
\usepackage[shortlabels]{enumitem}
\usepackage{float}
\newtheorem{definition}{Definition}
\newtheorem{proposition}{Proposition}
\newtheorem{theorem}{Theorem}
\newcommand{\defeq}{\vcentcolon=}
\renewcommand{\geq}{\geqslant}
\DeclareMathOperator{\im}{\operatorname{im}}
\renewcommand{\leq}{\leqslant}
\title{Finite group schemes as fundamental group schemes of smooth projective varieties}
\author{Gabriel Bassan}
\date{Source: arXiv:2608.27246v1 (CC BY 4.0)}
\begin{document}
\maketitle

\noindent\emph{Source and licence.} Finite group schemes as fundamental group schemes of smooth projective varieties. Version arXiv:2608.27246v1 (CC BY 4.0), distributed by arXiv under the Creative Commons Attribution 4.0 International licence, http://creativecommons.org/licenses/by/4.0/, as recorded in the arXiv licence metadata for this exact version and shown on the abstract page https://arxiv.org/abs/2608.27246v1. Full licence notice: \texttt{licenses/arxiv-2608.27246-CC-BY-4.0.txt}.\par\smallskip

\noindent\textit{Editorial scope.} Counted unit: the article's theorem Bertini (source0.tex lines 448--459). The article's complete original proof of it is delivered with it (source0.tex lines 460--499, one \texttt{proof} environment of 40 lines). No mathematics is written, repaired or re-derived: the statement and the proof are the article's own lines, in the article's own notation, and no retained line is substituted or re-flowed.  Retained uncounted as the source-local context the proof needs: lines 326--344; lines 346--357; lines 390--400; lines 402--412.  Nothing was omitted from any retained span, so the package carries no omission record.  No retained line contains a citation, so the wrapper carries no bibliography and no bibliography entry.  No retained line contains a figure environment, an image-inclusion command or drawing code, so no image is delivered and the package declares no asset dependency.  Editorial formatting only, in addition: an AMS \texttt{amsart} preamble that restates the article's own declarations and macros used by the retained lines, namely the environments \texttt{definition}, \texttt{proposition}, \texttt{theorem} on separate counters, together with the standard packages those lines need.  The retained environments print in this document as Definition 1, Proposition 1, Theorem 1 in order of appearance, whereas the article prints the same items with the numbering of its own full document.  The complete native source of the article as distributed in the arXiv e-print for this exact version is bundled unchanged as \texttt{sources/208700/source0.tex}.
\begin{definition}{(The universal section and zero locus)}\label{universal locus def}
Let $X/k$ be a finite type scheme of dimension $n$ and $1\leq c\leq n$ an integer. Let $L$ be a line bundle over $X$, $V$ be a finite dimensional $k$-vector space and $\varphi\colon V\to H^{0}(X,L)$ be a $k$-linear map such that its image generates $L$. We define the universal sections
    \[
    \sigma_{1},\dots, \sigma_{c}\in H^{0}(X\times V^{c}, p_{1}^{\ast}L)
    \]
    of $V$ as the image of the tautological $i$-th component sections 
    \[
    \tau_{1},\dots,\tau_{c}\in H^{0}(X\times V^{c}, V\otimes \mathcal{O}_{X\times V^{c}})
    \]
    by the morphism
    \[
    V\otimes \mathcal{O}_{X\times V^{c}}\to p_{1}^{\ast}L
    \]
    induced by $\varphi$. We also define the universal zero locus $\mathcal{Z}$ of $V$ as
    \[
    \mathcal{Z}\defeq V(\sigma_{1},\dots, \sigma_{c})\subseteq X\times V^{c}.
    \]
    We denote the maps $\mathcal{Z}\to X$ and $\mathcal{Z}\to V^{c}$ by $p_{X}$ and $p_{V}$.
\end{definition}

\begin{proposition}\label{universal locus prop}
Under the hypothesis of Definition \ref{universal locus def}, we have the following.
    \begin{enumerate}
        \item $\mathcal{Z}$ is isomorphic as an $X$-scheme to $\mathbb{V}_{X}(\ker(V\otimes\mathcal{O}_{X}\to L)^{c})$. In particular, it is a vector bundle of rank $c(\dim(V)-1)$ over $X$;
        \item The codimension of $\mathcal{Z}$ as a closed subscheme of $X\times V^{c}$ is equal to $c$;
        \item For $v=(v_{1},\dots, v_{c})\in V^{c}$, the fiber of the induced $k$-morphism $p_{V}\colon \mathcal{Z}\to V^{c}$ over $v$ is isomorphic to $V(\varphi(v_{1}),\dots, \varphi(v_{c}))\subseteq X$;
        \item If $p_{V}\colon \mathcal{Z}\to V^{c}$ is dominant, then for a generic $v\in V^{c}$ we have
        \[
        \dim(V(\varphi(v_{1}),\dots, \varphi(v_{c})))=n-c.
        \]
    \end{enumerate}
\end{proposition}

\begin{definition}{(Universal singular locus)}\label{singular locus def}
    Let $X/k$ be a smooth scheme of dimension $n$ and $1\leq c\leq n$ an integer. Let $L$ be a line bundle over $X$, $V$ be a finite dimensional $k$-vector space and $\varphi\colon V\to H^{0}(X,L)$ be a $k$-linear map such that its image generates $L$. Let $\mathcal{Z}$ be the universal zero locus and consider
    \[
    d\sigma_{1}\dots d\sigma_{c}\in H^{0}(\mathcal{Z},i_{\mathcal{Z}}^{\ast}(\Omega^{1}_{X\times V/V}(p_{X}^{\ast}L))).
    \]
    We define the universal singular locus $\mathcal{S}\subseteq \mathcal{Z}$ as
    \[
    \mathcal{S}\defeq V(d\sigma_{1}\wedge\cdots\wedge d\sigma_{c})\subseteq \mathcal{Z}.
    \]
    We will denote the induced morphisms $\mathcal{S}\to X$ and $\mathcal{S}\to V^{c}$ also by $p_{X}$ and $p_{V}$.
\end{definition}

\begin{proposition}\label{singular locus prop}
Under the hypothesis of Definition \ref{singular locus def}, we have the following.
    \begin{enumerate}
        \item Let $v=(v_{1},\dots, v_{c})\in V^{c}$. If $V(\varphi(v_{1}),\dots, \varphi(v_{c}))$ has codimension $c$, then the fiber $S_{v}$ over $v$ is isomorphic to the singular locus of $V(\varphi(v_{1}),\dots, \varphi(v_{c}))$.
        \item Suppose further that the induced $\mathcal{O}_{X}$-linear map
        \[
        p^{1}_{V}\colon V\otimes \mathcal{O}_{X}\to P^{1}_{X/k}(L)
        \]
        given by $v\otimes f\mapsto f\cdot P^{1}(\varphi(v))$ has constant rank $r>c$. Then the codimension of $\mathcal{S}$ as a closed subscheme of $\mathcal{Z}$ is equal to $r-c$.
    \end{enumerate}
\end{proposition}

\begin{theorem}\label{Bertini}
    Let $k$ be a field. Let $X/k$ be a smooth scheme over $k$ of pure dimension $n$ and $1\leq c\leq n$ be an integer. Let $L$ be a line bundle over $X$, $V$ be a $k$-vector space, $\varphi\colon V\to H^{0}(X,L)$ be a $k$-linear map and consider the $\mathcal{O}_{X}$-linear map
    \[
    p^{1}_{V}\colon V\otimes_{k}\mathcal{O}_{X}\to P^{1}(L).
    \]
    Suppose that
    \begin{enumerate}
        \item The image of $\varphi$ generates $L$;
        \item $p^{1}_{V}$ has constant rank $r>c$;
    \end{enumerate}
    Then for a generic $v\in V^{c}$, the zero locus $V(\varphi(v_{1}),\dots, \varphi(v_{c}))$ has dimension $n-c$ and its singular locus has dimension at most $n-r$.
\end{theorem}

\begin{proof}
    First of all, we can suppose that $k$ is algebraically closed since passing to the larger field extension does not affect dimensions or smoothness. In this case, the projection from the universal zero locus $p_{V}\colon \mathcal{Z}\to V^{c}$ is dominant. Indeed, it is enough to show that its image contains a nonempty open set. For this, let
    \[
    K\defeq \ker(\im p^{1}_{V}\to L)\subseteq \im p^{1}_{V}.
    \]
    Observe that for a fixed closed point $x\in X(k)$, since $\dim K(x)=r-1\geq c$, we can find $c$ linearly independent vectors $w_{1},\dots, w_{c}\in K(x)$. Since $V\to \im p^{1}_{V}$ is surjective, we can lift each $w_{i}$ to vectors $v_{i}\in V$ which satisfy 
    \begin{enumerate}
        \item $\varphi(v_{i})(x)=0$ for every $i$ and
        \item $d\varphi(v_{1})(x),\dots, d\varphi(v_{c})(x)$ are linearly independent in $\Omega^{1}_{X/k}(x)$.
    \end{enumerate}
    Let $v\defeq (v_{1},\dots, v_{c})$ Now, let $U$ be an open neighborhood of $x$ where $L$ is trivialized, let $e\in \Gamma(U\times V^{c},p_{1}^{\ast}L)$ be a generator and write $\sigma_{i}=f_{i}\cdot e$ with $f_{i}\in \Gamma(X\times V^{c},\mathcal{O}_{X\times V^{c}})$. Under the identification
    \[
    \Omega^{1}_{X\times V^{c}/V^{c}}\xrightarrow{\sim} p_{1}^{\ast}\Omega^{1}_{X/k}
    \]
    the vector $df_{i}(x,v)\in \Omega^{1}_{X\times V^{c}/V^{c}}(x,v)$ is sent to $ds_{i}(x)\in p_{1}^{\ast}\Omega^{1}_{X/k}$. Thus, the $df_{i}(x,v)$'s are linearly independent in $\Omega^{1}_{X\times V^{c}/V^{c}}(x,v)$ and, $p_{V}$ is smooth at $(x,v)$. Since a smooth morphism is open, we conclude that the image of $p_{V}$ contains a nonempty open set. In particular, by Proposition \ref{universal locus prop}, we have 
    \[
    \dim(\mathcal{Z}_{v})=n-c
    \]
    for a generic $v\in V^{c}$. By Proposition $\ref{singular locus prop}$, the fiber $\mathcal{S}_{v}$ of $\mathcal{S}$ over a point $v\in V^{c}$ such that $\mathcal{Z}_{v}$ has dimension $n-c$ is exactly its singular locus. We would like to calculate its dimension. For that, let $\{\mathcal{S}_{i}\}_{i\in I}$ be the (finitely many) irreducible components of $\mathcal{S}$. Let $I_{1}\subseteq I$ be the subset of indices such that 
    \[
    \overline{p_{V}(\mathcal{S}_{i})}=V^{c}
    \]
    and $I_{2}=I\setminus I_{1}$. For each $i\in I_{1}$, the morphism $p_{V,i}\colon \mathcal{S}_{i}\to V^{c}$ is dominant and therefore, the generic fiber satisfies
    \[
    \dim p_{V,i}^{-1}(v)=\dim \mathcal{S}_{i}-\dim V^{c}\leq \dim \mathcal{S}-\dim V^{c}=n-r.
    \]
    Let $U_{1}\subseteq V^{c}$ be a dense open subset for which the preceding inequality holds for every $i\in I_{1}$ and
    \[
    U_{2}=V^{c}\setminus(\bigcup_{i\in I_{2}} \overline{p_{V,i}(\mathcal{S}_{i})}).
    \]
    Then for every closed point $v\in U_{1}\cap U_{2}$ we have
    \[
    \dim p_{V}^{-1}(v)= \max_{i\in I}\dim p_{V,i}^{-1}(v)\leq n-r.
    \]
    Restricting ourselves to the dense open of $V$ where $\mathcal{Z}_{v}$ has dimension $n-c$, for a generic $s\in V^{c}$,
    \[
    \text{dim}(\text{Sing}(\mathcal{Z}_{v}))\leq n-r
    \]
    and we are done.
\end{proof}

\end{document}
